\chapter{Fermi liquids, Green's functions and response}
\setwlchapter{response}
\section{Resolvent or thermal trace?}
There are two ways to compute a retarded propagator on a finite
cluster. Commuting $A$ with $H$ generates a space of operators.
If it closes with $[O_i,H]=\sum_j M_{ij}O_j$, solving a linear
system at $z=\omega+i\eta$, $\eta>0$, sums all nested commutators.
Alternatively, a Lehmann sum uses the many-body energies and
transition matrix elements. The first route can remain symbolic;
the second is a numerical reference that includes correlations
without a factorization.
\casestudy{propagation}{How large an operator space does a chain require?}
\casestudy{lehmann}{Positive diagonal weight versus signed cross weight}
\section{High-frequency checks}
For an electron the anticommutator fixes the $1/z$ coefficient of
the diagonal propagator to one. The next moment probes the mean
addition energy, and a higher moment can distinguish a single pole
from a pair of poles with the same centroid. For a conserved density
the ordinary commutator with itself is zero: a nonzero thermal
fluctuation need not produce a nonzero dynamical retarded response.
\casestudy{moments}{An independent sum rule without finding the poles}

\section{A static limit is a prescription}
Define $\chi^R_{AB}(t)=-i\theta(t)\langle[A(t),B]\rangle$.
With a perturbation $+h B$, $\delta\langle A\rangle=\chi^R_{AB} h$
in this convention; the common Zeeman coupling is instead $-h S^z$.
For the density of a free band the retarded bubble (spin sum included)
is
\[
 \chi^0(q,z)=\frac{1}{N_s}\sum_{k\sigma}
 \frac{f(\epsilon_k)-f(\epsilon_{k+q})}{z+\epsilon_k-\epsilon_{k+q}} .
\]
At $q=0$ and finite $z$ the numerator vanishes. Taking the static
long-wavelength limit instead gives the negative thermodynamic density
of states for a perturbation raising the level energy. In the random
phase approximation a repulsive density interaction $v_q$ induces an
additional potential $v_q\delta n$, giving
$\chi=\chi^0/(1-v_q\chi^0)$. This argument fixes the denominator sign
without importing a convention for a positively defined polarization.
A current response to a vector potential also has a diamagnetic term:
the current itself depends on the field. Omitting it can produce a
spurious response to a pure gauge~\cite{vucicevic}.
\subsection{A two-level susceptibility by two independent routes}
For one electron of fixed spin on a dimer with hopping $-t$, the
bonding state has energy $-t$. The imbalance $D=n_1-n_2$ connects
it to the antibonding state at $+t$ with unit matrix element. Thus
$\chi^R_{DD}(z)=1/(z-2t)-1/(z+2t)$. Unlike total charge, this density
does not commute with the hopping. Check the static value at $t=1$
by applying the field directly. In the site basis $+hD$ changes the
diagonal energies in opposite directions:
\begin{wl}
dop = DiagonalMatrix[{1,-1}];
hd[h_] := {{h,-1},{-1,-h}};
{vals,vecs} = Eigensystem[hd[0]];
j = First[Ordering[vals,1]]; l = 3-j;
delta = vals[[l]]-vals[[j]];
wgt = Abs[Conjugate[Normalize[vecs[[j]]]].dop.Normalize[vecs[[l]]]]^2;
TableForm[{{delta,wgt,-2 wgt/delta,D[-Sqrt[1+h^2],{h,2}] /. h -> 0}},
 TableHeadings -> {None,{"Gap","Weight","Retarded limit","Curvature"}}]
\end{wl}
The last column is the second derivative of the exact lower eigenvalue
of $H+hD$. Its agreement with the retarded limit uses a nondegenerate
ground state and a probe independent of field. A thermal mixture
also changes its populations, as the exercise below illustrates.

\section{Quasiparticle parameters and available phase space}
Write $G^{-1}=z-\epsilon_k+\mu-\Sigma(k,z)$. For a normal Fermi
liquid at a regular point of its Fermi surface, the pole strength is
$Z=[1-\partial_\omega\operatorname{Re}\Sigma(k_F,0)]^{-1}$.
The Fermi velocity is $Z[\nabla_k\epsilon_k+
\nabla_k\operatorname{Re}\Sigma(k,0)]_{k_F}$ in units $\hbar=1$.
Only for a momentum-independent self-energy can the entire mass
enhancement of an isotropic parabolic band be read off from $1/Z$.
At first order a static Hartree shift changes the Fermi energy but
gives no decay. At second order a particle can excite a particle--hole
pair; thermal factors restrict all three final quasiparticles to a
shell about the Fermi surface. Two independent low-energy integrations
remain after energy conservation. With nonsingular phase space in
three dimensions the retarded imaginary part is negative and of order
$-\big[\omega^2+(\pi T)^2\big]$. In two dimensions logarithms can
modify this law. Broadening the discrete lines of a small cluster is
not a measurement of that lifetime.
\begin{exercise}[Fixed frequency or static field?]\label{ex:response}
For a single free level of energy $\epsilon$ at temperature $T>0$,
find the retarded density commutator and the isothermal derivative
$\partial\langle n\rangle/\partial\epsilon$ for one spin. Explain why
the two answers differ.
\end{exercise}
